% =====================================================================
%  Section: Frequency-domain (modal) analysis of the simplified drivetrain
%  Required in preamble:
%    \usepackage{amsmath, amssymb}
%    \usepackage{graphicx}
%    \usepackage{booktabs}
%    \usepackage{siunitx}   % (optional; used for \si/\SI/\num)
%  Adjust the figure path (fig/...) to your own.
% =====================================================================

\section{Frequency-Domain (Modal) Analysis of the Drivetrain}
\label{sec:frequency-analysis}

The purpose of this section is to characterise the linear oscillatory modes of
the simplified wind-turbine model, and in particular to locate and interpret
the mechanical \emph{torsional mode} of the shaft. The motivation is that this
mode is the candidate for an electrically driven mechanical resonance: the
converter (and hence the grid) can, in principle, pump energy into a lightly
damped torsional mode through the electromagnetic torque. The analysis proceeds
by linearising the full differential--algebraic model and then performing an
eigenvalue and eigenvector (modal) decomposition.

\subsection{Shaft model}
\label{subsec:drivetrain}

The drivetrain is modelled as a \emph{two-mass torsional system}: a heavy
rotor/hub mass with inertia constant $H_m$, coupled through a flexible shaft of
stiffness $K$ and damping $D$ to a much lighter generator mass $H_e$. The
drivetrain state vector is
\begin{equation}
  x_{\mathrm{dt}} =
  \begin{bmatrix} \omega_m & \omega_e & \theta_m & \theta_e \end{bmatrix}^{\!\top},
  \label{eq:dt-states}
\end{equation}
where $\omega_m,\omega_e$ are the rotor and generator speeds (pu) and
$\theta_m,\theta_e$ the corresponding angles. The equations of motion in
per-unit are
\begin{align}
  2H_m\,\dot{\omega}_m &= T_m - T_{\mathrm{sh}}, &
  2H_e\,\dot{\omega}_e &= T_{\mathrm{sh}} - T_e, \label{eq:swing}\\
  \dot{\theta}_m &= \omega_m, &
  \dot{\theta}_e &= \omega_e, \label{eq:angle}
\end{align}
with the shaft torque
\begin{equation}
  T_{\mathrm{sh}} = K_{pu}\,\theta_s + D_{pu}\,\omega_s,
  \qquad \theta_s = \theta_m-\theta_e, \qquad \omega_s = \omega_m-\omega_e.
  \label{eq:shaft}
\end{equation}
Here $T_m = P_{\mathrm{aero}}/\omega_m$ is the aerodynamic torque and
\begin{equation}
  T_e = \frac{P_e}{\omega_{e,\mathrm{filt}}\,\eta}
  \label{eq:Te}
\end{equation}
is the electromagnetic torque, where $P_e$ is the power the converter actually
delivers to the grid terminal and $\eta$ is the efficiency. Equation
\eqref{eq:Te} is the key coupling term: a disturbance in the delivered power
$P_e$ propagates directly into the generator-mass equation of motion
\eqref{eq:swing} and hence into the torsional dynamics.

\paragraph{Per-unit scaling.}
With $\omega_{\mathrm{base}} = \omega_{m,\mathrm{rated}}$ (in
\si{\radian\per\second}) and rated power $S_n$,
\begin{equation}
  T_{\mathrm{base}} = \frac{S_n}{\omega_{\mathrm{base}}}, \quad
  H = \frac{\tfrac12 J\,\omega_{\mathrm{base}}^2}{S_n}, \quad
  K_{pu} = \frac{K}{T_{\mathrm{base}}}, \quad
  D_{pu} = \frac{D\,\omega_{\mathrm{base}}}{T_{\mathrm{base}}}.
  \label{eq:pu}
\end{equation}
One convention detail is worth noting: the angle derivative in
\eqref{eq:angle} is written $\dot{\theta} = \omega$ without an
$\omega_{\mathrm{base}}$ factor. The consequence is that the modal frequency
comes out \emph{directly in physical} \si{\radian\per\second}, as confirmed
numerically below.

\subsection{Linearisation and the state matrix}
\label{subsec:linearisation}

The full model is a differential--algebraic system
\begin{equation}
  \dot{x} = f(x, v), \qquad 0 = g(x, v),
  \label{eq:dae}
\end{equation}
where $x$ collects the states and $v$ the bus voltages fixed by the algebraic
network equations $g(\cdot)=0$ (current injections against the network
admittance matrix). About an equilibrium point $(x_0, v_0)$, obtained from the
power flow and initialisation of the dynamic simulation, a first-order Taylor
expansion gives
\begin{equation}
  \Delta \dot{x} = A\,\Delta x,
  \qquad
  A = f_x - f_v\, g_v^{-1} g_x,
  \label{eq:Amatrix}
\end{equation}
where the sub-matrices are Jacobians evaluated at $(x_0,v_0)$, and the term
$-f_v g_v^{-1} g_x$ represents the elimination of the algebraic voltage
variables. The matrix $A$ is the reduced state matrix of the system.

\subsection{Eigenvalues: frequency and damping}
\label{subsec:eigenvalues}

The eigenvalue decomposition of $A$ is written
\begin{equation}
  A = R\,\Lambda\,L, \qquad
  \Lambda = \operatorname{diag}(\lambda_1,\dots,\lambda_n), \qquad
  L = R^{-1},
  \label{eq:eig}
\end{equation}
where the columns $r_i$ of $R$ are the right eigenvectors and the rows
$\ell_i^{\top}$ of $L$ are the left eigenvectors, normalised such that
$\ell_i^{\top} r_j = \delta_{ij}$. Each complex-conjugate eigenvalue
$\lambda_i = \sigma_i \pm j\omega_i$ corresponds to an oscillatory mode with
\begin{equation}
  f_i = \frac{\omega_i}{2\pi} \quad [\si{\hertz}],
  \qquad
  \zeta_i = \frac{-\sigma_i}{\sqrt{\sigma_i^2 + \omega_i^2}}
          = \frac{-\sigma_i}{|\lambda_i|},
  \qquad
  \tau_i = \frac{-1}{\sigma_i},
  \label{eq:fzeta}
\end{equation}
where $f_i$ is the damped oscillation frequency, $\zeta_i$ the damping ratio
and $\tau_i$ the time constant for the decay of the mode amplitude. A mode is
stable when $\sigma_i < 0$; the closer it lies to the imaginary axis
($\zeta_i \to 0$), the more lightly damped and resonance-prone it is.

\subsection{Mode shapes and participation factors}
\label{subsec:modeshape}

Two complementary quantities interpret a mode physically.

\paragraph{Mode shape (right eigenvector).}
The column $r_i$ gives the relative amplitude and phase of each state when
mode $i$ is excited alone:
\begin{equation}
  \Delta x(t) = \operatorname{Re}\!\left\{ r_i\, c\, e^{\lambda_i t} \right\},
  \label{eq:modeshape}
\end{equation}
for a scalar $c$ set by the initial condition. The mode shape thus shows
\emph{what} the mode looks like, including whether states oscillate in phase or
in anti-phase. We normalise so that the dominant component has unit length.

\paragraph{Participation factor.}
The dimensionless participation factor
\begin{equation}
  p_{ki} = \ell_{ik}\, r_{ki}
  \label{eq:pf}
\end{equation}
combines the observability ($r_{ki}$) and controllability ($\ell_{ik}$) of
state $k$ in mode $i$, and is invariant under scaling of the states. It
identifies which states dominate the dynamics of the mode, and is used to
classify the modes (mechanical, electromechanical, converter, etc.).

\subsection{Analytical characterisation of the torsional mode}
\label{subsec:torsion}

The system \eqref{eq:swing}--\eqref{eq:shaft} has two mechanical degrees of
freedom: a rigid common mode ($\lambda = 0$, the whole drivetrain rotating
together) and a torsional mode in which the two masses twist against each
other. Introduce the difference coordinates $\theta_s = \theta_m-\theta_e$ and
$\omega_s = \omega_m-\omega_e$. With no external torque,
\eqref{eq:swing}--\eqref{eq:angle} give
\begin{equation}
  \dot{\omega}_s
  = -\,T_{\mathrm{sh}}\!\left(\frac{1}{2H_m} + \frac{1}{2H_e}\right),
  \qquad \dot{\theta}_s = \omega_s .
  \label{eq:diff}
\end{equation}
Substituting \eqref{eq:shaft} yields a damped harmonic oscillator
\begin{equation}
  \ddot{\theta}_s + 2\zeta_{\mathrm{t}}\omega_{n}\,\dot{\theta}_s
  + \omega_{n}^{2}\,\theta_s = 0,
  \label{eq:osc}
\end{equation}
with natural frequency and damping ratio
\begin{equation}
  \boxed{\;\omega_{n} = \sqrt{K_{pu}\!\left(\frac{1}{2H_m}+\frac{1}{2H_e}\right)}\;},
  \qquad
  \zeta_{\mathrm{t}}
  = \frac{D_{pu}\!\left(\dfrac{1}{2H_m}+\dfrac{1}{2H_e}\right)}{2\,\omega_{n}}.
  \label{eq:wn}
\end{equation}

\paragraph{Mode shape, analytically.}
In the torsional mode the inertia-weighted common motion vanishes. Summing the
two equations in \eqref{eq:swing} (with no external torque) gives
$H_m\,A_m + H_e\,A_e = 0$, i.e.\ the amplitude ratio
\begin{equation}
  \frac{A_m}{A_e} = -\frac{H_e}{H_m} = -\frac{J_e}{J_m}.
  \label{eq:ratio}
\end{equation}
Because the generator mass is far lighter than the rotor ($J_e \ll J_m$), the
generator swings with much larger amplitude than the rotor, and the two are in
exact anti-phase. The mode is therefore \emph{generator-dominated}.

\subsection{Numerical results}
\label{subsec:results}

With the drivetrain parameters of the model
($S_n = \SI{15}{\mega\voltampere}$, $J_m = \num{3.525e8}$,
$J_e = \num{1.837e6}$, $K = \num{6.974e8}\,\si{\newton\metre\per\radian}$,
$D = \num{3.570e6}\,\si{\newton\metre\second\per\radian}$,
$\omega_{m,\mathrm{rated}} = \SI{7.56}{rpm}$), equation \eqref{eq:pu} gives
\begin{equation}
  \omega_{\mathrm{base}} = \SI{0.792}{\radian\per\second},\;
  H_m = \SI{7.36}{\second},\;
  H_e = \SI{0.0384}{\second},\;
  K_{pu} = 36.8,\;
  D_{pu} = 0.149 .
  \label{eq:num-pu}
\end{equation}
Substituted into \eqref{eq:wn}:
\begin{equation}
  \frac{1}{2H_m}+\frac{1}{2H_e} = 0.068 + 13.03 = 13.10,
\end{equation}
\begin{equation}
  \omega_n = \sqrt{36.8 \times 13.10} = \SI{21.95}{\radian\per\second}
  \;\Rightarrow\; f_n = \SI{3.49}{\hertz},
  \qquad
  \zeta_{\mathrm{t}} = \frac{0.149 \times 13.10}{2 \times 21.95} = 0.044 .
  \label{eq:num-result}
\end{equation}
The full modal analysis of the linearised model returns the torsional mode as
the eigenvalue
\begin{equation}
  \lambda_{\mathrm{t}} = -0.960 \pm j\,21.934
  \;\Rightarrow\;
  f_{\mathrm{t}} = \SI{3.491}{\hertz},\quad
  \zeta_{\mathrm{t}} = \SI{4.37}{\percent},\quad
  \tau_{\mathrm{t}} \approx \SI{1.0}{\second},
  \label{eq:num-eig}
\end{equation}
in excellent agreement with the analytical estimate \eqref{eq:num-result}. The
small discrepancy is due to coupling with the converter and network states,
which the isolated derivation of \S\ref{subsec:torsion} neglects.

The right eigenvector confirms the mode shape quantitatively. Normalised to the
dominant component,
\begin{equation}
  |r_{\omega_e}| = 0.998\;(\angle 180^\circ), \quad
  |r_{\omega_m}| = 0.0052\;(\angle 0^\circ), \quad
  |r_{\theta_e}| = 0.045, \quad
  |r_{\theta_m}| = 0.0002,
  \label{eq:num-rev}
\end{equation}
i.e.\ $|r_{\omega_m}|/|r_{\omega_e}| = 0.0052 = J_e/J_m$, exactly as predicted
by \eqref{eq:ratio}. The participation factors point to the same conclusion:
the torsional mode is completely dominated by the generator states $\omega_e$
and $\theta_e$ (both $\approx 1.0$), while the rotor contribution is negligible
($\approx 0.01$).

\begin{table}[t]
  \centering
  \caption{Selected oscillatory modes of the linearised model, sorted by
  frequency. The torsional mode is highlighted.}
  \label{tab:modes}
  \begin{tabular}{lccl}
    \toprule
    Mode & $f$ [\si{\hertz}] & $\zeta$ [\%] & Dominant states \\
    \midrule
    Rotor/MPT (slow)          & 0.064 & 23.1 & $\omega_m,\ \theta_m$ \\
    Speed LPF / pitch         & 0.109 & 78.1 & $\omega_{e,\mathrm{filt}}$, pitch \\
    Sync.\ gen.\ electromech.\ (COI) & 0.226 & 87.2 & gen.\ speed \\
    Sync.\ gen.\ $e_q'$       & 1.208 & 14.7 & $e_q'$ \\
    Converter (UIC) $v_{i,x}$ & 1.445 & 96.1 & $v_{i,x}$ \\
    Sync.\ gen.\ $e_q'$ (fast)& 3.228 & 28.9 & $e_q'$, gen.\ speed \\
    \textbf{Drivetrain torsional}& \textbf{3.491} & \textbf{4.37} &
      $\boldsymbol{\omega_e,\ \theta_e}$ \emph{(gen.-dominated)} \\
    Sync.\ gen.\ $e_d'$       & 5.204 & 0.26 & $e_d'$ \\
    Converter (UIC) $v_{i,y}$ & 7.998 & 25.9 & $v_{i,y},\ x_{\mathrm{filter}}$ \\
    \bottomrule
  \end{tabular}
\end{table}

Table~\ref{tab:modes} shows that the torsional mode at \SI{3.49}{\hertz} is the
only lightly damped \emph{mechanical} mode, and that it is well separated from
the neighbouring modes. Figure~\ref{fig:eigs} shows the eigenvalue spectrum in
the $s$-plane with the torsional mode highlighted, while
Figure~\ref{fig:modeshape} shows the mode shape in which the two masses swing
in anti-phase.

\begin{figure}[t]
  \centering
  % \includegraphics[width=0.8\linewidth]{fig/WT_LEOGO_eigs_full.png}
  \caption{Eigenvalue spectrum ($s$-plane) of the simplified model. The
  torsional mode ($\lambda = -0.96 \pm j21.93$, \SI{3.49}{\hertz},
  $\zeta = \SI{4.37}{\percent}$) is marked. The fact that the mode lies close
  to the imaginary axis is precisely why it can be driven to resonance.}
  \label{fig:eigs}
\end{figure}

\begin{figure}[t]
  \centering
  % \includegraphics[width=0.55\linewidth]{fig/WT_LEOGO_modeshape_torsional.png}
  \caption{Mode shape of the torsional mode. The generator ($\omega_e$) and the
  rotor ($\omega_m$) swing in exact anti-phase ($180^\circ$); the light
  generator mass carries almost all of the motion, with amplitude ratio
  $|r_{\omega_m}|/|r_{\omega_e}| = J_e/J_m \approx 0.0052$.}
  \label{fig:modeshape}
\end{figure}

\subsection{Resonance and electro-mechanical interaction}
\label{subsec:resonance}

The low damping gives a quality factor
\begin{equation}
  Q = \frac{1}{2\zeta_{\mathrm{t}}} \approx \frac{1}{2 \cdot 0.044} \approx 11,
  \label{eq:Q}
\end{equation}
i.e.\ a harmonic forcing at $f_n$ is amplified about \num{11} times relative to
a forcing far from resonance. The decisive point for the problem at hand is
\emph{where} the forcing can enter. The electromagnetic torque \eqref{eq:Te}
acts directly on the generator mass \eqref{eq:swing}, and because the torsional
mode is generator-dominated (large $|r_{\omega_e}|$ and participation) it is
both observable and controllable from the electrical port. A periodic power
disturbance
\begin{equation}
  P_e(t) = P_0 + \hat{P}\sin(2\pi f\, t)
  \label{eq:Pe-forcing}
\end{equation}
produces an oscillating $T_e$; as $f \to f_n$ the torsional response
$\theta_s,\omega_s,T_{\mathrm{sh}}$ builds up towards a steady-state amplitude
proportional to $Q$.

This is the linear analogue of \emph{sub-synchronous torsional interaction}
(SSTI/SSR): an electrical source---here the converter and the connected
grid---exchanging energy with a lightly damped mechanical torsional mode. In
the full aeroelastic FMU model this path is effectively closed, since the
torque controller (ROSCO) internally computes the generator torque from the
measured speed and thereby rejects external power/torque commands. In the
simplified model, by contrast, the path is \emph{open by construction} through
\eqref{eq:Te}, which makes it a suitable tool for demonstrating the phenomenon.
The experimental verification---injecting a \SI{3.49}{\hertz} electrical
disturbance and detecting the build-up in $T_{\mathrm{sh}}$, $\omega_s$ and
$\theta_s$, compared against an off-resonance case---is treated in
\S\ref{sec:forcedresponse}.
